\chapter{实数}\label{ch:b1:reals}

整个分析学都建立在一条性质之上，这条性质使 $\R$ 区别于
$\Q$：每个有\omterm{def:b1:reals:bounds}{上界}的非空\omterm{def:b1:logic:sets}{集合}都有\emph{最小}
\omterm{def:b1:reals:bounds}{上界}。本章精确陈述这一性质，并推出它的最初几个
推论——\omterm{thm:b1:reals:archimedes}{阿基米德性质}、\omterm{thm:b1:reals:floor}{取整函数}，以及有理数与无理数的
稠密性——并建立下文频繁使用的术语（sup、inf、max、min），
它们自
\cref{ch:b1:seq}起不断使用。

\section{上界性质}

\begin{definition}[界、上确界与下确界]\label{def:b1:reals:bounds}
设 $A \subseteq \R$ 为非空\omterm{def:b1:logic:sets}{集合}。实数 $M$ 为 $A$ 的\emph{上
界}\index{上界}，当 $a \leq M$ 对所有 $a \in A$ 成立时；
称 $A$ 为\emph{有上界}，指它存在上界（类似地
可讨论下界与有下界的情形；\emph{有界}指两者兼有）。
$A$ 的\emph{最大值}是属于 $A$ 的上界。

\emph{上确界}$\sup A$\index{上确界}是 $A$ 的最小上界，
若其存在；\emph{下确界}$\inf A$\index{下确界}是
最大下界。
\end{definition}

\begin{theorem}[$\R$ 的完备性公理]\label{thm:b1:reals:sup}
$\R$ 是包含 $\Q$ 的有序域，其中\emph{每个
有\omterm{def:b1:reals:bounds}{上界}的非空子集都有\omterm{def:b1:reals:bounds}{上确界}}。\index{
$\R$ 的完备性}
\end{theorem}
\admitted

\begin{remark}
将此作为 $\R$ 的定义性公理；构造模型（
可用戴德金分割或有理数的柯西序列）并证明其
唯一性是可行的，但篇幅较长，留待以后学习。注意
$\Q$ 不具有该性质：\omterm{def:b1:logic:sets}{集合} $\{x \in \Q : x^2 < 2\}$ 有\omterm{def:b1:reals:bounds}{上界}，
但没有\emph{在 $\Q$ 中}的最小\omterm{def:b1:reals:bounds}{上界}——其候选值
$\sqrt 2$ 缺失（见\cref{ex:b1:logic:sqrt2}）。通过
取相反数（$\sup(-A) = -\inf A$），每个有下界的非空\omterm{def:b1:logic:sets}{集合}
都有\omterm{def:b1:reals:bounds}{下确界}。
\end{remark}

\begin{proposition}[$\varepsilon$ 刻画]\label{prop:b1:reals:epsilon}
设 $A \neq \emptyset$ 有\omterm{def:b1:reals:bounds}{上界}，且 $s \in \R$。则 $s =
\sup A$ 当且仅当
\begin{enumerate}
  \item $s$ 为\omterm{def:b1:reals:bounds}{上界}：$\forall a \in A$，$a \leq s$；且
  \item 没有更小的数是\omterm{def:b1:reals:bounds}{上界}：$\forall \varepsilon > 0$，$\exists a \in
        A$, $a > s - \varepsilon$.
\end{enumerate}
\end{proposition}

\begin{proof}
若 $s = \sup A$，则（1）由定义即得；对于（2），$s -
\varepsilon < s$ 不是\omterm{def:b1:reals:bounds}{上界}，这恰好说明存在
$a > s - \varepsilon$。反之，（1）表明 $s$ 为\omterm{def:b1:reals:bounds}{上界}；
（2）表明任何 $t < s$ 都不是\omterm{def:b1:reals:bounds}{上界}（取 $\varepsilon
= s - t$），故 $s$ 为最小的\omterm{def:b1:reals:bounds}{上界}。
\end{proof}

\begin{example}\label{ex:b1:reals:supexamples}
$\sup \intoo{0}{1} = 1$，但不能取到（不存在最大值）；
$\sup \intcc{0}{1} = 1 = \max$。对于 $A = \{1 - \frac 1n : n \in
\N^*\}$：$\sup A = 1$，不能取到；$\inf A = \min A = 0$。最大值
（若存在）即为\omterm{def:b1:reals:bounds}{上确界}。引入 $\sup$ 的根本意义在于
在最大值不存在时提供替代。
\end{example}

\begin{omfigure}
\begin{tikzpicture}[scale=8.5]
  \draw[-Stealth] (-0.05,0) -- (1.3,0);
  \fill[omDef] (0,0) circle (0.35pt) node[below=2pt, font=\small]
    {$0$};
  \fill[omDef] (0.5,0) circle (0.35pt) node[below=2pt, font=\small]
    {$\tfrac12$};
  \fill[omDef] (0.6667,0) circle (0.35pt)
    node[below=2pt, font=\small] {$\tfrac23$};
  \fill[omDef] (0.75,0) circle (0.35pt)
    node[below=2pt, font=\small] {$\tfrac34$};
  \fill[omDef] (0.8,0) circle (0.35pt);
  \fill[omDef] (0.8333,0) circle (0.35pt);
  \fill[omDef] (0.8571,0) circle (0.35pt);
  \fill[omDef] (0.875,0) circle (0.35pt);
  \fill[omDef] (0.8889,0) circle (0.35pt);
  \fill[omDef] (0.9,0) circle (0.35pt);
  \draw[omThm, thick] (1,-0.015) -- (1,0.015);
  \node[above, omThm, font=\small] at (1,0.02) {$\sup A = 1$};
  \draw[omProp, very thick] (1,-0.04) -- (1.27,-0.04);
  \node[below, omProp, font=\small] at (1.13,-0.05)
    {$A$ 的上界};
  \draw[Stealth-, gray] (0.93,0.03) -- (0.8,0.075)
    node[above left, font=\small, text=gray, inner sep=1pt]
    {某个 $a > 1 - \varepsilon$};
\end{tikzpicture}

{\small \omterm{def:b1:logic:sets}{集合} $A = \{1 - \frac1n : n \in \N^*\}$ 在数轴上的图示：其中的点
向 $1$ 聚集，而不达到它。每个
$\geq 1$ 的数都是\omterm{def:b1:reals:bounds}{上界}（即那条射线），而更小的数都不是
\omterm{def:b1:reals:bounds}{上界}，因为 $A$ 中总有一个元素落入每个
\omterm{prop:b1:reals:intervals}{区间} $\intoo{1 - \varepsilon}{1}$ 之中，这正是
\cref{prop:b1:reals:epsilon} 的两个条件在同一幅图中的体现。\omterm{def:b1:reals:bounds}{上确界}是
\omterm{def:b1:reals:bounds}{上界}射线的左端点——而完备性
公理恰恰保证这条射线总有左
端点。}
\end{omfigure}

\begin{example}[实际计算上确界]\label{ex:b1:reals:supcompute}
利用\cref{prop:b1:reals:epsilon}计算两个完整示例。

\emph{\omterm{def:b1:logic:sets}{集合} $A = \{x + \frac1x : x > 0\}$。}对于每个 $x > 0$，
$x + \frac1x - 2 = \frac{(\,\sqrt x - 1/\sqrt x\,)^2}{1} \geq 0$,
故 $2$ 为下界；又 $2 = 1 + \frac11 \in A$，因此
$\inf A = \min A = 2$，在 $x = 1$ 处取到。另一方面，$A$
无\omterm{def:b1:reals:bounds}{上界}（$x + \frac1x > x$ 可超过任意 $M$，由
\cref{thm:b1:reals:archimedes}），故 $\sup A$ 在 $\R$ 中不存在
（$+\infty$ 即其在 $\overline\R$ 中的值）。

\emph{\omterm{def:b1:logic:sets}{集合} $B = \bigl\{\frac{m}{m + n} : m, n \in
\N^*\bigr\}$。}每个元素都属于 $\intoo{0}{1}$，故 $0$ 与
$1$ 均为界。两者都取不到：$\frac{m}{m+n} = 1$ 将
迫使 $n = 0$。求\omterm{def:b1:reals:bounds}{上确界}时，固定 $n = 1$，让 $m$ 增大：
$\frac{m}{m+1} = 1 - \frac{1}{m+1} > 1 - \varepsilon$ 只要
$m + 1 > \frac1\varepsilon$（由\omterm{thm:b1:reals:archimedes}{阿基米德性质}）：$\sup B = 1$。
类似地（$m = 1$，$n$ 充分大），$\inf B = 0$。收尾的要点是：
要确定\omterm{def:b1:reals:bounds}{上确界}，\omterm{def:b1:logic:sets}{集合}内只需一条选得恰当的\emph{单参数
路径}——此处即 $n = 1$——而
$\varepsilon$ 刻画便不再要求更多。
\end{example}

\begin{example}[下确界的对偶]\label{ex:b1:reals:infmirror}
\omterm{def:b1:reals:bounds}{下确界}有其自身的 $\varepsilon$ 刻画，由
\cref{prop:b1:reals:epsilon} 借助 $\inf A = -\sup(-A)$ 得出：
$i = \inf A$ 当且仅当 $i$ 为 $A$ 的下界，且对每个
$\varepsilon > 0$，存在 $a \in A$ 使 $a < i + \varepsilon$。
下面的例子同时处理两个界，设
\[
A = \Bigl\{(-1)^n + \frac1n : n \in \N^*\Bigr\}
= \Bigl\{0,\ \tfrac32,\ -\tfrac23,\ \tfrac54,\ -\tfrac45,\
\dots\Bigr\} .
\]
偶数指标给出 $1 + \frac1n \leq \frac32$，等号在
$n = 2$ 时取等。由于奇数指标处的值也都 $\leq 0 <
\frac32$，故 $\sup A = \max A = \frac32$。奇数指标给出
$-1 + \frac1n > -1$，并递减趋于 $-1$：
$A$ 的每个元素都 $> -1$，而 $-1 + \varepsilon$ 之下仍有奇数指标的 $-1 +
\frac1n$，只要 $n > \frac1\varepsilon$：$\inf A = -1$，
不可取到。同一个\omterm{def:b1:logic:sets}{集合}展示了四种情形：一个\omterm{def:b1:reals:bounds}{上确界}恰是最大值，
一个\omterm{def:b1:reals:bounds}{下确界}却不是最小值。
\end{example}

\begin{remark}[上、下确界的常见错误]
四类错误造成了大部分失分。(i) \emph{混淆
$\sup$ 与 $\max$}：$\sup A$ 未必属于 $A$；
$\max$ 仅当已在 $A$ 中找到一个是\omterm{def:b1:reals:bounds}{上界}的元素时才可用。
(ii) \emph{对\omterm{def:b1:reals:bounds}{上确界}传递严格不等式}：若 $a <
b$ 对所有 $a \in A$ 成立，则只能得出 $\sup A \leq b$——
反例：$A = \intoo{0}{1}$，$b = 1$。(iii) \emph{在检验合法性前写出 $\sup A$
}：该记号要求 $A$ 非空且
有\omterm{def:b1:reals:bounds}{上界}（见\cref{met:b1:reals:supproofs}）；$\sup \emptyset$
与 $\sup \N$ 在 $\R$ 中均无定义（$\overline\R$ 中的约定
需另行明确给出）。(iv) \emph{\omterm{def:b1:logic:sets}{集合}运算
}：$\sup(A \cup B) = \max(\sup A, \sup B)$ 恒成立，但
对 $A \cap B$ 没有一般结论——它可能为空，并且
即便非空，$\sup(A \cap B)$ 也可能远小于
$\min(\sup A, \sup B)$。例如取 $A = \{0, 2\}$，$B = \{0, 3\}$，
此时 $\sup(A \cap B) = 0$。
\end{remark}

\begin{example}[有限集有最大值——一个被默默使用的引理]\label{ex:b1:reals:finitemax}
每个非空有限集 $F \subseteq \R$ 都有最大值（并且
也有最小值）。对元素个数作归纳：单点集
$\{a\}$ 有 $\max = a$；若结论对 $n$ 个元素的\omterm{def:b1:logic:sets}{集合}成立，
而 $F$ 有 $n + 1$ 个元素，任取 $a \in F$：记\omterm{def:b1:logic:sets}{集合} $F
\setminus \{a\}$ 的最大值为 $m$，则 $\max F$ 为 $m$（当 $a
\leq m$ 时），否则为 $a$。无需用到完备性——这纯粹是序加归纳，
在 $\Q$ 中即已成立——然而此引理值得明确陈述一次，因为后文证明
会默默使用它：下文的取整构造（“夹在有限范围内的一组整数
有最大元”）、
在\cref{ch:b1:seq}中每一处形如
$\max(\abs{u_0}, \dots, \abs{u_{N-1}}, \dots)$ 的界、
以及在\cref{ch:b1:continuity}中“从有限多个
$\delta$ 中取最大者”的做法。无穷\omterm{def:b1:logic:sets}{集合}正是最大值失效之处，
\omterm{def:b1:reals:bounds}{上确界}接替其位置：本章正是为
无穷情形而设。
\end{example}

\begin{theorem}[阿基米德性质]\label{thm:b1:reals:archimedes}
对每个 $x \in \R$，存在 $n \in \N$ 使 $n > x$。等价地：
对所有 $\varepsilon > 0$ 与 $y > 0$，存在某个倍数 $n\varepsilon$
超过 $y$。\index{阿基米德性质}
\end{theorem}

\begin{proof}
假设不然，则某个 $x$ 为 $\N$ 的\omterm{def:b1:reals:bounds}{上界}。于是 $s = \sup \N$
存在（见\cref{thm:b1:reals:sup}）。由
\cref{prop:b1:reals:epsilon}~(2) 取 $\varepsilon = 1$，存在 $n
\in \N$ 使 $n > s - 1$；但此时 $n + 1 \in \N$ 且 $n + 1 > s$，
与 $s$ 为\omterm{def:b1:reals:bounds}{上界}相矛盾。对第二种形式，
设 $\varepsilon > 0$ 与 $y > 0$：将第一种形式用于 $x
= \frac{y}{\varepsilon}$，得到 $n \in \N$ 使 $n >
\frac{y}{\varepsilon}$，两边乘以 $\varepsilon > 0$
（乘正数保持严格不等号），得 $n\varepsilon > y$。
反之，取 $\varepsilon = 1$、$y = x$ 的第二种形式
即得 $x > 0$ 时的第一种形式，而 $n = 1$ 处理 $x \leq 0$ 的情形：
两个陈述严格等价。
\end{proof}

\begin{example}[阿基米德性质的应用]\label{ex:b1:reals:archimedeswork}
三个直接的应用，后文将经常用到。(i) \emph{没有正实数
低于所有 $\frac1n$}：若 $0 < \varepsilon$，
取 $n > \frac1\varepsilon$，则 $\frac1n < \varepsilon$。换言之，
$\R$ 中不含无穷小——通常所说的
“$\frac1n$ 任意变小”这一非形式说法，正是本定理。
(ii) \emph{显式阈值}：$n$ 须多大，才使
$\frac{1}{n^2} \leq 10^{-6}$？只需 $n \geq 10^3$ 即可。
——\omterm{thm:b1:reals:archimedes}{阿基米德性质}保证此类 $n$ 存在，代数计算
则将其定出。(iii) \emph{幂次超过任意界}：$2^n \geq n +
1$（归纳可证），因此对每个 $M$，某个 $2$ 的幂超过 $M$：
这正是在\cref{exo:b1:reals:8}中处理二进有理数时所用的几何增长。
收尾的要点是：\omterm{thm:b1:reals:archimedes}{阿基米德
性质}是每句形如“取
$n$ 充分大”的说法背后的依据——此后我们可自由使用这一说法，
而本例即为其一次性的论证。
\end{example}

\begin{theorem}[取整函数]\label{thm:b1:reals:floor}
对每个 $x \in \R$，恰有一个整数，即
\emph{取整函数}\index{取整函数} $\lfloor x \rfloor$，满足
\[
\lfloor x \rfloor \leq x < \lfloor x \rfloor + 1 .
\]
\end{theorem}

\begin{proof}
\emph{存在性。}设\omterm{def:b1:logic:sets}{集合} $E = \{k \in \Z : k \leq x\}$ 非空：
由\cref{thm:b1:reals:archimedes}，存在 $m \in \N$，使 $m > -x$，
于是 $-m < x$，从而 $-m \in E$。该\omterm{def:b1:logic:sets}{集合}有\omterm{def:b1:reals:bounds}{上界}（任一
满足 $n > x$ 的整数均可作\omterm{def:b1:reals:bounds}{上界}，此类整数同理存在），因此
该\omterm{def:b1:logic:sets}{集合}是落在有限\omterm{prop:b1:reals:intervals}{区间} $\intint{-m}{n}$ 中的整数集，
故有最大元 $k = \max E$。于是 $k \leq x$，而 $k + 1 \notin E$
意味着 $x < k + 1$。

\emph{唯一性。}若 $k$ 与 $k'$ 均满足上述不等式，
则由 $k \leq x < k' + 1$ 得 $k \leq k'$，对称地有 $k'
\leq k$。
\end{proof}

\begin{example}[取整函数的应用]\label{ex:b1:reals:floorpractice}
$\lfloor 3.7 \rfloor = 3$，$\lfloor 5 \rfloor = 5$，且 $\lfloor
-3.7 \rfloor = -4$：取整是\emph{向下}取值，而非朝 $0$ 靠拢。
由\cref{thm:b1:reals:floor}中的唯一性，可得以下两个结论，
我们将不加说明地使用：其一，对 $n \in \Z$，
\[
\lfloor x + n \rfloor = \lfloor x \rfloor + n ,
\]
因为 $\lfloor x \rfloor + n$ 是整数，且满足
$x + n$ 的两个定义不等式——而满足它们的整数只有一个。
其二，$\lfloor \, \cdot \, \rfloor$ 是非减的：若 $x \leq
y$，则 $\lfloor x \rfloor \leq x \leq y < \lfloor y \rfloor + 1$，
而整数若 $< \lfloor y \rfloor + 1$，则必 $\leq \lfloor y
\rfloor$。但需注意，一般地 $\lfloor 2x \rfloor \neq 2\lfloor
x \rfloor$：取 $x = 0.6$，则 $\lfloor 1.2 \rfloor = 1
\neq 0 = 2\lfloor 0.6 \rfloor$。

真正\emph{成立}的是一个值得记住的恒等式（埃尔米特恒等式，
其最简单的情形）：对每个实数 $x$，
\[
\lfloor x \rfloor + \Bigl\lfloor x + \frac12 \Bigr\rfloor
= \lfloor 2x \rfloor .
\]
记 $x = \lfloor x\rfloor + u$，其中 $u \in \intco{0}{1}$，
分两种情形。若 $u < \frac12$，则左边为 $\lfloor
x\rfloor + \lfloor x\rfloor = 2\lfloor x\rfloor$，而 $2x =
2\lfloor x\rfloor + 2u$，且 $2u \in \intco{0}{1}$，故右边同样为
$2\lfloor x\rfloor$。若 $u \geq \frac12$，则左边为
$\lfloor x\rfloor + (\lfloor x\rfloor + 1)$，又 $2u \in
\intco{1}{2}$，故右边为 $2\lfloor x\rfloor + 1$。
最后值得注意的是：$\lfloor x + \frac12\rfloor$ 是 $x$ 的
\emph{四舍五入}（取最近整数），因此该恒等式表明：
即“取整与四舍五入之和，等于二倍之数的取整”——
而对小数部分 $u$ 作分情形讨论，
是处理各类取整恒等式的标准技巧。
（\cref{exo:b1:reals:2,exo:b1:reals:3} 亦依此法。）
\end{example}

\begin{theorem}[$\Q$ 与 $\R \setminus \Q$ 的稠密性]\label{thm:b1:reals:density}
对任意两个实数 $x < y$，其间必有一个有理数
与一个无理数。\index{稠密性}
\end{theorem}

\begin{proof}
\emph{存在有理数。}由\cref{thm:b1:reals:archimedes}，取 $n \in
\N^*$，使 $n > \frac{1}{y - x}$，从而 $ny - nx > 1$。令 $m = \lfloor
nx \rfloor + 1$。一方面，$nx < \lfloor nx \rfloor + 1 = m$
（由\cref{thm:b1:reals:floor}）；另一方面，$m = \lfloor nx \rfloor +
1 \leq nx + 1 < ny$。两边除以 $n$，得 $x < \frac mn < y$。

\emph{存在无理数。}对数对 $x -
\sqrt 2 < y - \sqrt 2$ 应用前一结论：存在有理数 $q$ 介于二者之间，于是
$q + \sqrt 2 \in \intoo{x}{y}$ 为无理数（若 $q + \sqrt 2$ 为有理数，
则 $\sqrt 2$ 亦为有理数）。
\end{proof}

\begin{example}[稠密性证明的实际运行]\label{ex:b1:reals:densityrun}
此证明即一个算法，下面就 $x = 1.414$ 与
$y = \sqrt 2$ 实际执行一遍。由于 $1.4142^2 = 1.99996164 < 2$，故 $\sqrt 2
> 1.4142$，从而 $y - x > 0.0002$，且 $\frac{1}{y - x} < 5000$：
取 $n = 5000$ 是合理的。于是 $nx = 7070$，故 $m = \lfloor
7070 \rfloor + 1 = 7071$，所得有理数为
\[
\frac{m}{n} = \frac{7071}{5000} = 1.4142,
\qquad
1.414 < 1.4142 < \sqrt 2 .
\]
值得注意的是：证明只需 $n$ 略大于
$\frac{1}{y-x}$，所得即为 $\frac 1n$ 的整数倍中
首个大于 $x$ 的数。稠密性并非抽象的奇迹，
而是伪装的长除法，这一主题在周末问题中有详尽展开
（\cref{pb:b1:reals:1}）。
\end{example}

\begin{remark}[完备性接下来在何处使用]
\Cref{thm:b1:reals:sup} 是本书唯一的非代数公理，
而分析中每一个存在性定理，都不过是该公理
换了装束：单调收敛定理
（\cref{ch:b1:seq}）、波尔扎诺--魏尔斯特拉斯定理
（\cref{ch:b1:topology}）、介值定理与极值
定理（\cref{ch:b1:continuity}），以及积分本身的定义——
即下和的\omterm{def:b1:reals:bounds}{上确界}（\cref{ch:b1:integration}）。
本科三年级分册即在同一公理之上建立测度论与希尔伯特空间。
此后各章中，若证明凭空得出一个实数，
请去寻找其中隐藏的\omterm{def:b1:reals:bounds}{上确界}。
\end{remark}

\begin{remark}[离散性与稠密性之间]
$\Z$ 与 $\Q$ 位于 $\R$ 中的两端：每个整数的周围
都有一段长为 $1$ 的间隙，其中不含别的整数
（离散性——正是它使\omterm{thm:b1:reals:floor}{取整函数}得以良定义），
而任意两实数之间有无穷多个有理数
（稠密性）。值得注意的是，对 $\R$ 的\emph{加法\omterm{def:b1:structures:subgroup}{子群}}
则不存在中间情形：\cref{exo:b1:reals:9} 证明了
此类\omterm{def:b1:structures:subgroup}{子群}要么具有 $\alpha\Z$ 的形式（离散），要么
稠密——这一二分法为 $\{\sin n\}$ 在
\cref{ch:b1:seq} 中的稠密性，以及
\cref{pb:b1:continuity:1} 中的构造性“怪物”提供依据。一般\omterm{def:b1:logic:sets}{集合}当然可以
自由地混合这两种行为：$\Z \cup \Q\cap\intcc{0}{1}$ 在远处是离散的，
而在中部则是稠密的。
\end{remark}

\begin{method}[用上确界与下确界证明等式]\label{met:b1:reals:supproofs}
欲证 $\sup A = s$，先验证 $s$ 为 $A$ 的\omterm{def:b1:reals:bounds}{上界}，再
对每个 $\varepsilon > 0$（或取 $\varepsilon = \frac
1n$），找出 $A$ 中大于 $s - \varepsilon$ 的元素。比较\omterm{def:b1:reals:bounds}{上确界}时，
可用：$A \subseteq B \implies \sup A \leq \sup B$；并对所有 $a, b$：
$\sup(A + B) = \sup A + \sup B$，其中 $A + B = \{a + b\}$
（见\cref{exo:b1:reals:5}）。切勿写出 $\sup A$，
除非已确知 $A$ 非空且有\omterm{def:b1:reals:bounds}{上界}。
\end{method}

\section{区间}

\begin{proposition}[区间的刻画]\label{prop:b1:reals:intervals}
子集 $I \subseteq \R$ 是\emph{区间}\index{区间}（
常见类型之一：$\intoo{a}{b}$、$\intcc{a}{b}$、$\intco{a}{b}$、
$\intoc{a}{b}$、半直线、$\R$、$\emptyset$、单点集），当且
仅当它是\emph{凸}的：
\[
\forall x, y \in I,\ \forall z \in \R, \quad
x \leq z \leq y \implies z \in I .
\]
\end{proposition}

\begin{proof}
所列各类型显然都是凸的。反之，设 $I$ 为凸集
且非空。取 $a = \inf I$（若 $I$ 有下界），否则取 $a =
-\infty$；同理，取 $b = \sup I$，否则取 $+\infty$。我们断言 $\intoo{a}{b}
\subseteq I \subseteq \intcc{a}{b}$（约定在
$\pm\infty$ 处取显然的含义）。第二个包含即为界的定义。
对于第一个包含，任取 $z \in \intoo{a}{b}$：由 $z > a$ 知 $z$ 不是
下界（若 $a = -\infty$ 则无此要求），故存在 $x \in I$ 使 $x < z$；
同理存在 $y \in I$ 使 $y > z$；由凸性即得 $z \in I$。

余下只需由双包含关系
$\intoo{a}{b} \subseteq I \subseteq \intcc{a}{b}$ 读出类型：
夹在开\omterm{prop:b1:reals:intervals}{区间}与其闭包之间的\omterm{def:b1:logic:sets}{集合}，与
$\intoo{a}{b}$ 的区别仅在于（有限）端点是否属于该\omterm{def:b1:logic:sets}{集合}。具体地，
若 $a, b \in \R$，则 $(a \in I,\ b \in I)$ 的四种可能
分别给出 $\intoo{a}{b}$、$\intco{a}{b}$、
$\intoc{a}{b}$、$\intcc{a}{b}$（包括退化情形
$a = b$：若 $a \in I$ 则为单点集）；若 $a = -\infty$ 且 $b \in
\R$，则得 $\intoo{-\infty}{b}$ 或 $\intoc{-\infty}{b}$；
对 $a \in \R$、$b = +\infty$ 对称地处理；$a = -\infty$、
$b = +\infty$ 时 $I = \R$。所有情形均在上述列表之中，证毕。
\end{proof}

\begin{remark}[为何凸性是正确的判别准则]
本命题将一个\emph{几何}式的定义（十种形状的列表）
化为一条简洁的\emph{逻辑}判别条件，而人们实际使用的
正是它：要证明某\omterm{def:b1:logic:sets}{集合}是\omterm{prop:b1:reals:intervals}{区间}，
切勿去追问它属于十种形状中的哪一种——只需验证凸性，
其类型交由该命题判定。\cref{ch:b1:continuity} 中的介值
定理将恰以此方式陈述（“\omterm{prop:b1:reals:intervals}{区间}的连续像
仍为\omterm{prop:b1:reals:intervals}{区间}”），其证明得到的
是凸性，而非
\omterm{prop:b1:reals:intervals}{区间}的具体形状。
\end{remark}

\begin{remark}[扩充实直线]
为方便起见，添加两个符号，并在 $\overline\R = \R
\cup \{-\infty, +\infty\}$ 中讨论，约定 $\sup A = +\infty$
当 $A$ 无\omterm{def:b1:reals:bounds}{上界}时，$\sup \emptyset = -\infty$。于是
$\R$ 的\emph{每一个}子集都在 $\overline\R$ 中有\omterm{def:b1:reals:bounds}{上确界}——这是一种
记号上的便利，在\cref{ch:b1:seq} 中讨论极限时被自由使用。
\end{remark}

\begin{example}[在 $\overline\R$ 中计算]\label{ex:b1:reals:extendedcompute}
依上述约定：$\sup \Z = +\infty$，$\inf \Z =
-\infty$；对 $A = \{n + (-1)^n n : n \in \N\} = \{0, 4, 0, 8,
\dots\} \cup \{0\}$，则 $\sup A = +\infty$（偶数项 $2n$
无界），$\inf A = \min A = 0$；且 $\sup\emptyset =
-\infty \leq \inf\emptyset = +\infty$——这是唯一一个
\omterm{def:b1:reals:bounds}{上确界}\emph{小于}\omterm{def:b1:reals:bounds}{下确界}的\omterm{def:b1:logic:sets}{集合}，提醒我们这些约定的选取是为了使
$\sup$ 关于包含单调递增、$\inf$
关于包含单调递减：
\[
A \subseteq B \implies \sup A \leq \sup B
\quad\text{且}\quad \inf A \geq \inf B ,
\]
此时无需非空条件即可成立。约定所
\emph{不}提供的是算术：$+\infty + (-\infty)$ 与 $0
\times (+\infty)$ 仍无定义，而对\omterm{def:b1:reals:bounds}{上确界}所作的每一步代数
变形，都须先检验是否会出现这些式子。
扩充实直线仅是记号工具，并非数系。
\end{example}

\begin{example}[逸出 $\Q$ 的上确界]\label{ex:b1:reals:supescape}
回到开篇注记中的\omterm{def:b1:logic:sets}{集合} $A = \{x \in \Q : x^2 <
2\}$，并求其\emph{在 $\R$ 中}的\omterm{def:b1:reals:bounds}{上确界}。它非空
（$1 \in A$），且以 $1.5$ 为\omterm{def:b1:reals:bounds}{上界}（若 $x > 1.5$，则 $x^2 >
2.25 > 2$），故 $s = \sup A$ 存在。我们断言 $s = \sqrt 2$（即
由\cref{exo:b1:reals:12} 构造的实数）。\omterm{def:b1:reals:bounds}{上界}：每个
$a \in A$ 满足 $a < \sqrt2$——当 $a \leq 0$ 时
显然成立，而当 $a > 0$ 时，若 $a \geq \sqrt2$ 则 $a^2 \geq 2$。
没有更小的\omterm{def:b1:reals:bounds}{上界}：给定 $t < \sqrt2$，由稠密性
（\cref{thm:b1:reals:density}）可取有理数 $q$ 使
$\max(1, t) < q < \sqrt 2$，于是 $q^2 < 2$，故 $q \in A$
大于 $t$。由\cref{prop:b1:reals:epsilon}，$s = \sqrt2
\notin \Q$。本例的要点在于：有理数\omterm{def:b1:logic:sets}{集合}的\omterm{def:b1:reals:bounds}{上确界}
未必是有理数——完备性恰恰是
$\R$（与 $\Q$ 不同）的承诺：决不让\omterm{def:b1:reals:bounds}{上确界}逸出；
本例即本章开篇的注记，如今已经证明，
而不再只是指出。
\end{example}

\begin{remark}[本分册内的展望]
本章三个工具今后各有其用途。
\omterm{def:b1:reals:bounds}{上确界}主导分析的一面：单调极限
（\cref{ch:b1:seq}）、积分的定义本身，
（\cref{ch:b1:integration}），以及在
\cref{ch:b1:euclid} 的几何中，点到子空间的距离——
一个\omterm{def:b1:reals:bounds}{下确界}，经正交投影化为最小值。\omterm{thm:b1:reals:floor}{取整
函数}在离散与连续相遇之处
反复出现：数字展开（本章的周末问题），
狄利克雷的抽屉原理逼近
（\cref{pb:b1:derivative:1}），和式的积分比较
（\cref{ch:b1:series}）。稠密性论证可升格为一种方法，
在\cref{ch:b1:continuity} 中：连续函数的恒等式
只需在 $\Q$ 上验证即可——柯西
函数方程的一半（\cref{pb:b1:continuity:1}）正是这一手法。
遇到疑问时，若想知道本分册某一证明的存在性
论断从何而来，答案几乎总是：本章。
\end{remark}

\section{练习}

\begin{exercise}[$\star$]\label{exo:b1:reals:1}
求下列\omterm{def:b1:logic:sets}{集合}的 sup、inf、max、min——若存在——并给出证明：
\[
A = \Bigl\{\frac{1}{n} : n \in \N^*\Bigr\},
\qquad
B = \Bigl\{\frac{(-1)^n n}{n+1} : n \in \N\Bigr\},
\qquad
C = \{x \in \R : x^2 < 3\}.
\]
\end{exercise}

\begin{exercise}[$\star$]\label{exo:b1:reals:2}
证明：对所有 $x, y \in \R$，$\lfloor x \rfloor + \lfloor y
\rfloor \leq \lfloor x + y \rfloor \leq \lfloor x \rfloor + \lfloor y
\rfloor + 1$，且两个界都可以取到。
\end{exercise}

\begin{exercise}[$\star$]\label{exo:b1:reals:3}
证明：对每个 $x \in \R$ 及 $n \in \N^*$，有
$\Bigl\lfloor \frac{\lfloor nx \rfloor}{n} \Bigr\rfloor = \lfloor x
\rfloor$。
\end{exercise}

\begin{exercise}[$\star$]\label{exo:b1:reals:4}
设 $A \subseteq B$ 为 $\R$ 的非空子集，且 $B$ 有界。证明
$\inf B \leq \inf A \leq \sup A \leq \sup B$。
\end{exercise}

\begin{exercise}[$\star\star$]\label{exo:b1:reals:5}
对非空有界的 $A, B \subseteq \R$，定义 $A + B = \{a + b : a
\in A,\ b \in B\}$，且 $-A = \{-a : a \in A\}$。证明：
\[
\sup(A + B) = \sup A + \sup B,
\qquad
\sup(-A) = -\inf A .
\]
\end{exercise}

\begin{exercise}[$\star\star$]\label{exo:b1:reals:6}
设 $f, g \colon E \to \R$ 为有界函数。证明
\[
\sup_{x \in E}\, \bigl(f(x) + g(x)\bigr) \leq \sup_{x \in E} f(x) +
\sup_{x \in E} g(x),
\]
并给出使不等式严格成立的例子。为何这并不与
\cref{exo:b1:reals:5} 相矛盾？
\end{exercise}

\begin{exercise}[$\star\star$]\label{exo:b1:reals:7}
证明 $\sqrt 2 + \sqrt 3$ 是无理数。\emph{（将其平方，并利用
$\sqrt 6$ 的无理性，其证明须借助
\cref{exo:b1:arith:7}。）}
\end{exercise}

\begin{exercise}[$\star\star$]\label{exo:b1:reals:8}
证明二进有理数集 $D = \bigl\{\frac{m}{2^n} : m \in \Z,\ n \in
\N\bigr\}$ 在 $\R$ 中稠密：任意两个
实数之间必有一个二进有理数。
\end{exercise}

\begin{exercise}[$\star\star\star$]\label{exo:b1:reals:9}
设 $G$ 为 $(\R, +)$ 的\omterm{def:b1:structures:subgroup}{子群}，且 $G \neq \{0\}$。记 $\alpha =
\inf\,(G \cap \intoo{0}{+\infty})$。证明：
\begin{enumerate}
  \item 若 $\alpha > 0$，则 $G = \alpha\Z$；
  \item 若 $\alpha = 0$，则 $G$ 在 $\R$ 中稠密。
\end{enumerate}
由此推出 $\Z + \sqrt 2\,\Z$ 在 $\R$ 中稠密。
\end{exercise}

\begin{exercise}[$\star\star\star$]\label{exo:b1:reals:10}
设 $A, B$ 为正实数的非空\omterm{def:b1:logic:sets}{集合}，令 $AB = \{ab : a \in
A, b \in B\}$。证明 $\sup(AB) = \sup A \cdot \sup B$（有界情形），
并举例说明正性条件不可或缺。
\end{exercise}

\begin{exercise}[$\star\star$]\label{exo:b1:reals:11}
对非空有界的 $A \subseteq \R$，定义
\emph{直径}
\[
\operatorname{diam} A = \sup\,\{\abs{a - a'} : a, a' \in A\} .
\]
证明 $\operatorname{diam} A = \sup A - \inf A$，且
$\intcc{\inf A}{\sup A}$ 是
包含 $A$ 的最小闭\omterm{prop:b1:reals:intervals}{区间}。
\end{exercise}

\begin{exercise}[$\star\star\star$]\label{exo:b1:reals:12}
设 $y > 0$，且 $E = \{x \geq 0 : x^2 \leq y\}$。证明 $E$ 是
非空的且有\omterm{def:b1:reals:bounds}{上界}，并证明 $s = \sup E$ 满足 $s^2 =
y$ \emph{（分别排除 $s^2 < y$ 与 $s^2 > y$，
在各情形中构造一个小的 $h > 0$，与
\omterm{def:b1:reals:bounds}{上确界}的定义相矛盾）}。由此推出每个 $y > 0$ 都有唯一的
正平方根 $\sqrt y > 0$，且 $y \mapsto \sqrt y$ 在
$\intoo{0}{+\infty}$ 上递增。
\end{exercise}

\section{周末问题：数字展开与有理数的节律}

\begin{problem}[{周末问题——$b$ 进制展开：
存在性、唯一性与周期性刻画 $\Q$}]
\label{pb:b1:reals:1}
$\intco{0}{1}$ 中的每个实数，在每个底
$b \geq 2$ 下都有数字展开；排除末尾
为数字 $b - 1$ 的串，展开唯一；且展开最终周期
当且仅当该数为有理数。本题证明上述三点，
只依据完备性公理——不用数列，也不用级数，
仅用\omterm{def:b1:reals:bounds}{上确界}、\omterm{thm:b1:reals:archimedes}{阿基米德性质}与\omterm{thm:b1:reals:floor}{取整函数}
——并以数字形式的康托对角线论证作结。
全文设 $b \geq 2$ 为固定整数，称为\emph{底}，
\emph{数字}是 $\intint{0}{b-1}$ 中的元素，且数字
串 $(d_n)_{n \geq 1}$ 称为\emph{正常的}，当其
不最终等于 $b - 1$（即：对每个 $N$，存在 $n >
N$，使 $d_n \leq b - 2$）。\index{$b$ 进制展开}
\index{小数展开}

\medskip
\noindent\textbf{第一部分——手算数位。}$p$ 除以
$q$ 以 $b$ 为底的长除法：将当前余数乘以 $b$，
除以 $q$，把商记为下一个数位，并保留余数。
\begin{enumerate}
  \item 以 $10$ 为底，对 $\frac 18$ 与
        $\frac 17$ 运行该算法，每一步记录数位\emph{和}余数。
        检验 $\frac 17$ 的余数依次循环取 $1, 3, 2, 6, 4, 5$，
        并检验数位 $142857$ 之后
        无限重复。
  \item 计算以 $2$ 为底的 $\frac 13$ 与
        $\frac{5}{16}$ 的展开，以及以 $3$ 为底的 $\frac 12$ 的展开。
        观察：一个数有限终止，另外两个数无限循环——
        而 $\frac 12$ 以 $10$ 为底有限终止，
        以 $3$ 为底则无限重复。
  \item 设既约分数 $x = \frac pq \in \intco{0}{1}$，证明：
        算法产生的数位最终全为 $0$，当且仅当
        余数 $b^N p \bmod q$ 对某个 $N$ 为零，当且仅当
        $q$ \omterm{def:b1:arith:divides}{整除}某个幂 $b^N$，当且仅当
        $q$ 的每个素因数都\omterm{def:b1:arith:divides}{整除} $b$。
        验证：$\frac{1}{20}$ 以 $10$ 为底有限终止，
        以 $3$ 为底则不然。
  \item 定义\emph{截断} $s_n = \lfloor b^n x \rfloor
        / b^n$。取 $x = \sqrt 2$，$b = 10$，计算 $s_0,
        \dots, s_4$，每一步通过验证两个相邻的平方数夹住 $2$
        （例如 $1.4142^2
        = 1.99996164 < 2 < 2.00024449 = 1.4143^2$），并检验
        $s_n \leq \sqrt 2 < s_n + 10^{-n}$ 在每一步都成立。
\end{enumerate}

\medskip
\noindent\textbf{第二部分——由\omterm{def:b1:reals:bounds}{上确界}得存在性。}固定
$x \in \intco{0}{1}$，记 $A_n = \lfloor b^n x \rfloor$ 与
$d_n = A_n - b\,A_{n-1}$，对 $n \geq 1$。
\begin{enumerate}[resume]
  \item 证明 $A_0 = 0$，且 $b\,A_{n-1} \leq A_n \leq b\,A_{n-1} +
        b - 1$；由此推出每个 $d_n$ 都是数位。
  \item 证明 $s_n := A_n b^{-n}$ 满足
        \[
        s_n = \sum_{k=1}^{n} d_k\,b^{-k}
        \qquad\text{且}\qquad
        s_n \leq x < s_n + b^{-n} .
        \]
  \item 用归纳法证明 $b^n \geq n + 1$，然后证明
        $(s_n)$ 单调不减，且 $x = \sup_n s_n$
        \emph{（利用 \cref{prop:b1:reals:epsilon} 与
        \cref{thm:b1:reals:archimedes}）}。
  \item 证明 $(d_n)$ 为真数位串：若 $d_k = b - 1$
        对所有 $k > N$ 成立，则以有限等比求和计算 $s_n$（$n > N$），
        并与第 6 题矛盾。
  \item 反之，设 $(e_n)_{n \geq 1}$ 为任一真数位串，并令
        $t_n = \sum_{k=1}^n e_k b^{-k}$。证明 $y
        = \sup_n t_n$ 存在，属于 $\intco{0}{1}$，并且
        满足 $t_n \leq y < t_n + b^{-n}$，对每个 $n$ 成立
        \emph{（严格不等式需取某个 $e_m \leq b
        - 2$ 且 $m > n$）}。由此得 $\lfloor b^n y \rfloor = b^n
        t_n$，进而 $y$ 的数位（按第 5 题的意义）
        恰为 $e_n$。
\end{enumerate}

\medskip
\noindent\textbf{第三部分——唯一性、序与平移。}
\begin{enumerate}[resume]
  \item 综合第 5--9 题，得到\emph{$b$ 进制展开
        定理}：\omterm{def:b1:logic:map}{映射} $x \mapsto (d_n)$ 与 $(e_n) \mapsto
        \sup_n t_n$ 是 $\intco{0}{1}$ 与真数位串\omterm{def:b1:logic:sets}{集合}之间
        互逆的\omterm{def:b1:logic:inj}{双射}。特别地，两个不同的真数位串
        不可能有相同的
        值。
  \item 现允许非真数位串。证明：若 $e_n
        = b - 1$ 对所有 $n > M$ 成立（$M \geq 0$ 取最小），则其值为
        $t_M + b^{-M}$；由此得 $0.999\dots = 1$ 以
        $10$ 为底成立，并且恰有两种数位表示的实数
        正是 $b$ 进制分数
        $m/b^N \in \intoo{0}{1}$——其余每个实数，即使计入非真数位串，
        也只有一种表示。
  \item 证明第 10 题中的\omterm{def:b1:logic:inj}{双射}对字典序保序：若
        $x$ 与 $y$ 的真数位串首次在
        下标 $m$ 处不同，则
        $x < y$ 当且仅当 $d_m < e_m$。
  \item （平移引理）设 $x \in \intco{0}{1}$ 的数位为
        $(d_n)$。证明 $bx$ 的小数部分的数位为
        $(d_{n+1})_{n \geq 1}$ \emph{（计算 $\lfloor b^n(bx -
        A_1)\rfloor$，利用 $\lfloor u - K \rfloor = \lfloor u
        \rfloor - K$，其中 $K$ 为整数）}，并用归纳法推出
        $b^m x$ 的小数部分的数位为
        $(d_{n+m})_{n \geq 1}$。
\end{enumerate}

\medskip
\noindent\textbf{第四部分——有理性即周期性。}设 $x
= \frac pq \in \intco{0}{1}$ 为既约分数，且 $r_n = b^n p
\bmod q$ 的余数，即 $b^n p$ 除以
$q$.
\begin{enumerate}[resume]
  \item 证明 $A_n = \dfrac{b^n p - r_n}{q}$ 与 $r_n =
        (b\,r_{n-1}) \bmod q$。
  \item 证明 $d_n = \Bigl\lfloor \dfrac{b\,r_{n-1}}{q}
        \Bigr\rfloor$：每个数位仅由前一个余数决定，
        这恰是第一部分中的
        长除法。
  \item 应用抽屉原理
        （见\cref{cor:b1:counting:pigeonhole}）于 $r_0, \dots, r_q$，
        推出：每个有理数的展开最终都是周期的，
        其前周期与周期都不超过
        $q$。
  \item 反之，设 $y \in \intco{0}{1}$ 的数位是
        \emph{纯}周期的：$d_{n+T} = d_n$ 对所有 $n
        \geq 1$ 成立。利用平移引理及
        第 10 题的唯一性，证明 $b^T y$ 的小数部分
        等于 $y$，于是推出 $(b^T - 1)\,y \in \N$，因此 $y$ 是
        有理数，其分母\omterm{def:b1:arith:divides}{整除} $b^T - 1$。在 $0.(142857)$ 上验证
        此机制：$142857 \times 7 = 999999$。
  \item 对最终周期的情形用平移处理，并叙述
        \emph{周期性判别准则}：$x \in \intco{0}{1}$ 为有理数，
        当且仅当其真 $b$ 进制展开
        最终是周期的——在某一进制下成立，
        当且仅当在所有进制下成立。
  \item 对 $x = \frac 1q$，$\gcd(q, b) = 1$，证明
        其展开为纯周期的，且最小周期为
        最小的 $T \geq 1$ 满足 $b^T \equiv 1 \pmod q$（即
        $b$ 模 $q$ 的乘法阶）。验证：当 $q = 7$，$b = 10$ 时，
        $10$ 的幂模 $7$ 依次为 $3, 2, 6, 4, 5, 1$，
        阶为 $6$，与
        第 1 题吻合。
\end{enumerate}

\medskip
\noindent\textbf{第五部分——红利与对角线。}
\begin{enumerate}[resume]
  \item 设 $x^*$ 为 $\intco{0}{1}$ 中这样的实数，其 $10$ 进制数字为 $1$，
        位于三角数位置 $\frac{j(j +
        1)}{2}$（$j \geq 1$）处，其余为 $0$，即 $x^* =
        0.101001000100001\dots$。证明其数字串正规而非最终周期
        \emph{（若周期为 $T$
        则 1 之间的间隔至多为 $T$，而间隔不断增大）}，
        由此可得 $x^*$ 为无理数：这是一个
        仅凭数字节奏便证明其无理的数。
  \item 证明对任意底 $b$，\omterm{def:b1:logic:sets}{集合} $\{m/b^n : m \in \Z,
        n \in \N\}$ 在 $\R$ 中稠密（推广
        \cref{exo:b1:reals:8}），且每个有理数 $\frac pq
        \in \intoo{0}{1}$ 在底 $q$ 下都有\emph{有限}展开。
        要旨：“有限”是（数，底）这一对的性质，
        而周期性——即有理性——则是
        数本身所固有的。
  \item （康托对角线法）设 $k \mapsto x_k$ 是从
        $\N^*$ 到 $\intco{0}{1}$ 的任一\omterm{def:b1:logic:map}{映射}。定义数字串 $e_k =
        1$，若第 $k$ 位数字（取自 $x_k$）不同于 $1$，
        否则取 $e_k = 2$。证明 $(e_k)$ 正规，
        其值 $y$ 属于 $\intco{0}{1}$，且 $y \neq
        x_k$ 对所有 $k$ 成立。由此得：不存在从 $\N^* \to
        \intco{0}{1}$ 的\omterm{def:b1:logic:inj}{满射}。（有关可数性的术语
        与本定理的恰当归属，
        见\cref{ch:b1:topology}。）
  \item 证明：若 $x$ 与 $y$ 的正规展开直至指标 $n$ 都一致，
        则 $\abs{x - y} < b^{-n}$，但逆命题不成立。取
        $x = 0.1$、$y = 0.0999$（底 $10$）为反例，
        可见数值接近并不保证数字一致。
        哪些实数应为此负责？
  \item 对 $x = \frac{1}{10}$，在底 $b = 2$ 下运行第四部分的算法：
        依次计算余数与数字，直至出现循环，并
        由此得 $\frac{1}{10} = (0.0\overline{0011})_2$，
        前期长 $1$，周期长 $4$。利用第 3 题说明为何
        任何有限二进制字符串都永远不会等于 $\frac{1}{10}$
——这正是计算机浮点数中 $0.1 + 0.2$ 并不恰好
        等于 $0.3$ 的原因。
  \item 综合。每项用一句话回答：证明在何处用到了
        (i) 完备性，(ii) \omterm{thm:b1:reals:archimedes}{阿基米德性质}，(iii)
        \omterm{thm:b1:reals:floor}{取整函数}的唯一性条款，(iv) 鸽巢
        原理？要点是：$\intco{0}{1}$ 可由正规数字串忠实编码，
        有理性可由周期性读出，
        然而分析更偏爱\omterm{def:b1:reals:bounds}{上确界}而非数字。为什么？
        （想一想两个数字串如何相加。）
\end{enumerate}
\end{problem}
